% Generated by tools/textbook/export_listings.js from tools/decks/mastery.js and tools/concepts/registry.json: do not edit.
\begin{chapstatus}{Stato di apprendimento di questo capitolo}
\masterychip{now}{DA COMPRENDERE OGGI}\enspace registri, calcolo, trasferimento, salto; i registri; contratto init/update/draw; gli argomenti in a0 ... a7; perché si salva ra; indirizzo, valore, load, store; le istruzioni di base; le pseudo-istruzioni; un'istruzione è un numero; seguire a mano l'esecuzione.\par\smallskip
\masterychip{use}{DA UTILIZZARE OGGI}\enspace la convenzione di chiamata (Capitolo~\ref{ch:processor}); le maschere di bit (Capitolo~\ref{ch:boole1}) e, usati senza essere studiati nel corso, le direttive dell'assemblatore, gli strumenti di PRG32, le chiamate del runtime PRG32, il debugger di PRG32-QT, il programma completo.\par\smallskip
\masterychip{later}{DA COMPRENDERE PIÙ AVANTI}\enspace la pila (Capitolo~\ref{ch:processor}); i modi di indirizzamento (Capitolo~\ref{ch:instructions}); formati delle istruzioni (Capitolo~\ref{ch:instructions}); l'esadecimale (Capitolo~\ref{ch:information}); il complemento a due (Capitolo~\ref{ch:integers}); la virgola fissa (Capitolo~\ref{ch:reals}); l'algebra di Boole (Capitolo~\ref{ch:boole1}); il sommatore (Capitolo~\ref{ch:boole2}); dai volt ai bit (Capitolo~\ref{ch:electronics}); latch e flip-flop (Capitolo~\ref{ch:sequential}); il ciclo del processore (Capitolo~\ref{ch:instructions}); il progetto Verm1 (Capitolo~\ref{ch:verm1p1}).
\end{chapstatus}
